% Part: second-order-logic
% Chapter: metatheory
% Section: loewenheim-skolem

\documentclass[../../../include/open-logic-section]{subfiles}

\begin{document}

\olfileid{sol}{met}{lst}

\olsection{દ્વિતીય-ક્રમ તર્કશાસ્ત્રમાં લેવેનહાઇમ--સ્કોલેમ પ્રમેય નિષ્ફળ જાય છે}

\begin{explain}
(અધોગામી) લેવેનહાઇમ--સ્કોલેમ પ્રમેય કહે છે કે !!{sentence}sના જે
ગણને અનંત નિદર્શ હોય, તેને !!a{enumerable} નિદર્શ પણ હોય છે. આ
પ્રમેય પણ પૂર્ણતા પ્રમેયનું પરિણામ છે: પૂર્ણતાની સાબિતી
!!{sentence}sના કોઈપણ સુસંગત ગણ માટે નિદર્શ રચે છે, અને તે નિદર્શ
!!{enumerable} હોય છે. ઊર્ધ્વગામી લેવેનહાઇમ--સ્કોલેમ પ્રમેય પણ છે;
તે ખાતરી આપે છે કે !!{sentence}sના ગણને !!a{denumerable} નિદર્શ
હોય તો તેને !!a{nonenumerable} નિદર્શ પણ હોય છે. દ્વિતીય-ક્રમ
તર્કશાસ્ત્રમાં બંને પ્રમેયો નિષ્ફળ જાય છે.
\end{explain}


\begin{thm}
\ollabel{thm:sol-no-ls} દ્વિતીય-ક્રમ તર્કશાસ્ત્રમાં લેવેનહાઇમ--સ્કોલેમ
પ્રમેય નિષ્ફળ જાય છે: એવા !!{sentence}s છે જેમને અનંત નિદર્શો છે,
પણ કોઈ !!{enumerable} નિદર્શ નથી.
\end{thm}

\begin{proof}
યાદ કરો કે
\[
\fn{Count} \ident \lexists[z][\lexists[u][\lforall[X][((X(z) \land
      \lforall[x][(X(x) \lif X(u(x)))]) \lif \lforall[x][X(x)])]]]
\]
!!a{structure}~$\Struct{M}$માં ત્યારે અને માત્ર ત્યારે જ સાચું છે
જ્યારે~$\Domain{M}$ !!{enumerable} હોય. તેથી
$\lnot\fn{Count}$ એ~$\Struct{M}$માં ત્યારે અને માત્ર ત્યારે જ
સાચું છે જ્યારે~$\Domain{M}$ !!{nonenumerable} હોય. આવી
!!{structure}s છે: કોઈપણ !!{nonenumerable} ગણને !!{domain} લો,
દાખલા તરીકે $\Pow{\Nat}$ અથવા~$\Real$. આથી
$\lnot \fn{Count}$ને અનંત નિદર્શો છે, પરંતુ કોઈ
!!{enumerable} નિદર્શ નથી.
\end{proof}

\begin{thm}
એવા !!{sentence}s છે જેમને !!{denumerable} નિદર્શો છે,
પણ કોઈ !!{nonenumerable} નિદર્શ નથી.
\end{thm}

\begin{proof}
$\fn{Count} \land \fn{Inf}$ એ~$\Nat$માં સાચું છે, પરંતુ
જે !!{structure}~$\Struct{M}$ માટે~$\Domain{M}$
!!{nonenumerable} હોય તેમાં સાચું નથી.
\end{proof}


\end{document}
